%! Linear Transformations and Matrices — Unit 4, Lesson 4.12 — Exit Ticket
\documentclass[10pt]{article}
\usepackage{precalculus-article}
\usepackage{precalculus-boxes}
\newcommand{\TallMath}[1]{$\displaystyle #1\rule[-1.4em]{0pt}{3.2em}$}

\begin{document}
\pageheader{Unit 4, Lesson 4.12}{Exit Ticket}
\namedateperiod
\vspace{0.12in}

\begin{notesbox}
\textbf{1.} Let $A = \begin{bmatrix}3&-1\\0&2\end{bmatrix}$ and
$\vec{v}=\begin{bmatrix}2\\1\end{bmatrix}$.
Compute $L(\vec{v})=A\vec{v}$. Show your work.

\writelines{5}

\vspace{0.04in}
$L(\vec{v}) = \begin{bmatrix}\hspace{1.8cm}\\[8pt]\hspace{1.8cm}\end{bmatrix}$
\end{notesbox}

\vspace{0.12in}

\begin{notesbox}
\textbf{2.} A linear transformation $L:\mathbb{R}^2\to\mathbb{R}^2$ maps
$\hat{\imath}=\langle1,0\rangle$ to $\langle-1,0\rangle$ and
$\hat{\jmath}=\langle0,1\rangle$ to $\langle0,-1\rangle$.
Write the $2\times2$ matrix $A$ associated with $L$.

\vspace{0.16in}
\hspace{1em}$A = \begin{bmatrix}\hspace{1.4cm}&\hspace{1.4cm}\\[10pt]\hspace{1.4cm}&\hspace{1.4cm}\end{bmatrix}$
\vspace{0.14in}
\end{notesbox}

\vspace{0.12in}

\begin{notesbox}
\textbf{3.} \textbf{True or False:} A linear transformation can map $\vec{0}$
to $\langle1,1\rangle$. Circle one: \quad \textbf{True} \qquad \textbf{False}

\vspace{0.08in}
Explain in one sentence: \writeline
\end{notesbox}

\end{document}
